% ============================================================================
% Fichier  : partie2/P02_C06_balayage.tex
% Titre    : Balayage et Énumération
% Auteur   : MINKA MI NGUIDJOÏ Thierry Emmanuel
% Version  : 1.0 -- novembre 2026
% ============================================================================
% Sources : RFC 9293 (TCP) ; NIST SP 800-115 (découverte) ; documentation de
% l'outil Nmap. Résultats de l'Opération LOGONE : cas fictif.
% ============================================================================

\chapter{Balayage et Énumération}
\label{chap:balayage}
\chaptermark{Balayage et énumération}

\epigraph{%
    « Un port ouvert est une question~:\\
    qui répond, et pourquoi est-il là~? »%
}{Principe du cours}

\niveauChapitre{Intermédiaire}{Chapitre~\ref{chap:osint}~; cours SEC3266,
chapitres sur les réseaux.}

\begin{boxobjectifs}
\begin{itemize}
    \item Expliquer la découverte d'hôtes et le balayage de ports TCP et UDP.
    \item Interpréter les états d'un port~: ouvert, fermé, filtré.
    \item Identifier les services et leurs versions, et énumérer leurs
          informations.
    \item Conduire un balayage maîtrisé, sans perturber la production.
    \item Relier le balayage à sa détection par le défenseur.
\end{itemize}
\end{boxobjectifs}

\bigskip

% ============================================================================
\section*{Question directrice : combien de portes ouvertes comptait vraiment LiZbiZ~?}
% ============================================================================

Le 21~octobre, M.~TAMO balaye la plage externe de LiZbiZ~; le 26, depuis le
poste fourni, Mme~NGO~MBOCK balaye les réseaux de Garoua. Les résultats
mettent en regard ce qui devrait être ouvert, selon la matrice des flux du
cours SEC3266, et ce qui l'est.

\medskip
\textbf{Plan du chapitre.}
\begin{enumerate}
    \item Découvrir les hôtes (section~\ref{sec:c06-decouverte}).
    \item Balayer les ports (section~\ref{sec:c06-ports}).
    \item Identifier les services (section~\ref{sec:c06-services}).
    \item Balayer sans nuire (section~\ref{sec:c06-maitrise}).
    \item Ce que voit le défenseur (section~\ref{sec:c06-defenseur}).
\end{enumerate}

\separateur

% ============================================================================
\section{Découvrir les Hôtes}
\label{sec:c06-decouverte}
% ============================================================================

La découverte cherche les adresses qui répondent~: requêtes ICMP, requêtes
ARP sur le réseau local, tentatives de connexion TCP vers des ports
courants. Un hôte qui ne répond pas au « ping » n'est pas forcément éteint~:
un pare-feu peut filtrer l'ICMP.

\begin{boxdef}
\termecle{Énumération}~: étape qui, après la découverte, extrait d'un
service les informations qu'il livre (noms de partages, comptes, versions,
bannières, configuration).
\end{boxdef}

% ============================================================================
\section{Balayer les Ports}
\label{sec:c06-ports}
% ============================================================================

TCP établit une connexion en trois temps (SYN, SYN-ACK, ACK) défini par la
RFC~9293~\cite{rfc9293}. Le balayage exploite les réponses à un SYN.

\begin{tableaucours}{L{2.4cm}L{4.6cm}L{5.2cm}}{%
    \tableauHeader{État} &
    \tableauHeader{Réponse observée} &
    \tableauHeader{Signification}
}
Ouvert & SYN-ACK & Un service écoute \\
Fermé & RST & L'hôte répond, aucun service \\
Filtré & Aucune réponse ou message ICMP d'interdiction & Un filtrage
intercepte \\
\end{tableaucours}
\vspace{-0.8em}
\begin{center}{\small\itshape Tableau~6.1. Les états d'un port TCP}\end{center}

Le balayage UDP est plus lent et plus incertain~: un service UDP ouvert ne
répond souvent rien, ce qui le confond avec un port filtré. On le réserve
aux ports utiles (DNS, SNMP, NTP).

% ============================================================================
\section{Identifier les Services}
\label{sec:c06-services}
% ============================================================================

\begin{boxscenario}{Opération LOGONE. Prévu et constaté}
\textbf{Constatations.} Externe~: 7~hôtes actifs sur 32~adresses, 18~ports
ouverts, dont 4 non prévus par la matrice des flux (dont l'ancienne
passerelle, constat T-03). Interne, Garoua~: 214~hôtes actifs~; le
protocole de partage de fichiers SMBv1 est actif sur deux serveurs~; une
imprimante expose son interface d'administration avec le mot de passe par
défaut (constat déjà signalé en mars, cours SEC3266).

\medskip
\textbf{Analyse} (confiance élevée). Les écarts entre la matrice et la
réalité sont autant de constats~: un service non prévu n'est ni surveillé,
ni corrigé.
\end{boxscenario}

\begin{tableaucours}{L{2.0cm}L{3.0cm}L{7.2cm}}{%
    \tableauHeader{Port} &
    \tableauHeader{Service} &
    \tableauHeader{Information à énumérer}
}
22 & SSH & Version, méthodes d'authentification acceptées \\
80, 443 & HTTP, HTTPS & Serveur, application, certificat, pages \\
445 & SMB & Version du protocole, partages accessibles, signature \\
389 & LDAP & Annuaire~: domaine, possibilité d'interrogation anonyme \\
161 (UDP) & SNMP & Communauté par défaut, informations système \\
\end{tableaucours}
\vspace{-0.8em}
\begin{center}{\small\itshape Tableau~6.2. Ports courants et énumération}\end{center}

% ============================================================================
\section{Balayer sans Nuire}
\label{sec:c06-maitrise}
% ============================================================================

\begin{boxprocedure}
\textbf{Un balayage maîtrisé}
\begin{enumerate}
    \item Vérifier le périmètre et les exclusions (Chapitre~\ref{chap:engagement}).
    \item Commencer par la découverte, puis les ports courants, puis le reste.
    \item Limiter le débit sur les réseaux fragiles (automates, imprimantes,
          liaisons lentes).
    \item Éviter les scripts intrusifs sur la production sans accord.
    \item Enregistrer les résultats bruts, horodatés, dans le journal.
\end{enumerate}
\end{boxprocedure}

\begin{boxalert}
\textbf{Violation typique.} Lancer un balayage complet à plein débit sur
une liaison inter-sites lente.

\smallskip
\textbf{Règle~:} un balayage peut saturer un équipement ou une liaison et
interrompre un service~: c'est la conséquence que réprime l'article~66 si
elle n'est pas maîtrisée. On adapte le débit au réseau.
\end{boxalert}

% ============================================================================
\section{Ce que Voit le Défenseur}
\label{sec:c06-defenseur}
% ============================================================================

Un balayage est bruyant~: une source qui contacte de nombreux ports ou de
nombreux hôtes en peu de temps est un motif classique (cours SEC3266). Le
SOC de LiZbiZ a détecté le balayage interne en 40~minutes.

\begin{boxmetier}{Analyste SOC et équipe rouge}
\textbf{Corréler les balayages internes.} Une règle simple~: alerte si une
même source interne contacte plus de 100~ports ou 50~hôtes en dix minutes.
Elle aurait vu PC-GAR-07 en mars (SEC3266) et le testeur en octobre.
\end{boxmetier}

\separateur

\begin{boxresume}
\textbf{Ce qu'il faut retenir de ce chapitre~:}
\begin{itemize}
    \item \textbf{Découverte, balayage, énumération}~: trois étapes.
    \item \textbf{Ouvert, fermé, filtré}~: trois états d'un port TCP~; UDP est
          plus incertain.
    \item \textbf{Comparer la réalité à la matrice des flux}~: chaque écart
          est un constat.
    \item \textbf{Balayer sans nuire}~: périmètre, progressivité, débit,
          scripts prudents.
    \item \textbf{Le balayage est bruyant}~: le SOC doit le voir.
\end{itemize}
\end{boxresume}

\begin{boxinfo}
\textbf{Points de contrôle de l'élève.} Je sais~: expliquer la découverte
d'hôtes~; interpréter l'état d'un port~; choisir quoi énumérer selon le
service~; conduire un balayage maîtrisé~; écrire une règle de détection
simple.
\end{boxinfo}

\bigskip

\begin{boxexo}[title={\ding{45}\quad Exercice~6.1. Les états \niveaubase}]
Quel état attribuer à un port qui répond par RST~? par SYN-ACK~? qui ne
répond pas~?
\end{boxexo}

\begin{boxexo}[title={\ding{45}\quad Exercice~6.2. Le ping \niveaubase}]
Pourquoi un hôte qui ne répond pas au « ping » peut-il être actif~?
\end{boxexo}

\begin{boxexo}[title={\ding{45}\quad Exercice~6.3. Les écarts \niveaubase}]
Sur les 18~ports ouverts à l'externe, combien étaient prévus par la
matrice des flux~?
\end{boxexo}

\begin{boxexo}[title={\ding{45}\quad Exercice~6.4. UDP \niveaumoyen}]
Pourquoi un balayage UDP est-il plus lent et plus incertain qu'un balayage
TCP~?
\end{boxexo}

\begin{boxexo}[title={\ding{45}\quad Exercice~6.5. SMBv1 \niveaumoyen}]
Rédigez la fiche de constat « SMBv1 actif sur deux serveurs ».
\end{boxexo}

\begin{boxexo}[title={\ding{45}\quad Exercice~6.6. L'imprimante \niveaumoyen}]
L'imprimante avait déjà été signalée en mars. Que révèle sa présence en
octobre sur la gestion des corrections~? (cours GES4062)
\end{boxexo}

\begin{boxexo}[title={\ding{45}\quad Exercice~6.7. Le débit \niveaumoyen}]
Proposez un plan de balayage des 214~hôtes de Garoua qui ne sature pas la
liaison avec Douala.
\end{boxexo}

\begin{boxexo}[title={\ding{45}\quad Exercice~6.8. Détecter \niveaumoyen}]
Le SOC a mis 40~minutes à voir le balayage. Proposez deux améliorations
pour réduire ce délai.
\end{boxexo}

\begin{boxexo}[title={\ding{45}\quad Exercice~6.9. Énumération \niveauexpert}]
Pour chacun des ports du Tableau~6.2, dites quelle mesure défensive
réduit l'information livrée à un attaquant.
\end{boxexo}

\begin{boxexo}[title={\ding{45}\quad Exercice~6.10. Inventaire continu \niveauexpert}]
Concevez un processus de balayage interne mensuel, autorisé, qui alimente
l'inventaire et la matrice des flux de LiZbiZ.
\end{boxexo}

\bigskip

\begin{boxinfo}
\textbf{Pour aller plus loin.}
\begin{itemize}
    \item RFC~9293 (TCP)~\cite{rfc9293}~; NIST SP~800-115~\cite{nist800115}.
\end{itemize}
\end{boxinfo}

\begin{boxinfo}
\textbf{Transition.} Un service identifié doit être confronté aux failles
connues~: le Chapitre~\ref{chap:vulnerabilites} apprend à les analyser et
à les prioriser.
\end{boxinfo}
